\clearpage

\chapter{Methodology}

In this chapter, we outline the methodology used to develop and evaluate the procedural
quest generation system for RPGs using LLMs. The goal of this research is to create an
automated framework that can generate coherent, engaging, and contextually appropriate
quests for RPGs, leveraging the power of LLMs and hybrid architectures. This chapter
describes the step-by-step procedure, from data collection and pre-processing to model
training, fine-tuning, and evaluation. The methodology is designed to ensure that the
generated quests are of high quality, both from a computational and a user experience
perspective.

We begin by discussing the dataset collection and preparation, including the specifics
of the quest data used to train the models. Next, we detail the architecture of the LLM-based
models, the fine-tuning procedures, and the use of LoRA adapters to optimize model
performance for quest generation. Following that, we delve into the evaluation metrics
employed to assess the generated quests, explaining both quantitative and qualitative
methods such as automatic evaluation metrics and human evaluation. Finally, we outline
the experimental setup and the procedures followed during model training, validation,
and testing, with a focus on how the methodology ensures the generation of diverse and
contextually rich quests for RPGs.

\section{Dataset Description and Preprocessing}

This section introduces the dataset compiled for training and evaluating LLMs in the task
of PCG. The construction of the dataset involves aggregating quest data from diverse
RPG sources, standardizing the structure into a uniform schema, and formatting it for
compatibility with generative language models.

Given the narrative and structural complexity of quests in RPGs, it is essential to
represent each quest in a machine-readable yet semantically rich format that preserves
critical gameplay elements such as objectives, subtasks, quest givers, rewards, and supporting
entities. This requires a preprocessing pipeline capable of handling heterogeneous
data from multiple games, some of which originate from structured datasets while others
are manually curated and normalized with the aid of LLMs.

The final dataset, referred to as the QUEST Dataset, captures a wide range of quest
archetypes and contextual variations across both high-fantasy and post-apocalyptic game
settings. It combines data from multiple titles, including games from the \textit{Elder Scrolls,
Fallout, and Baldur's Gate franchises}~\cite{theelderscrollsivoblivion,theelderscrollsvskyrim,fallout1,fallout2,baldursgate,baldursgate2shadowsofamn}, and ensures each entry adheres to a unified
tagging convention that facilitates token-level alignment during training.

To prepare the dataset for sequence modeling tasks, each quest instance is transformed
into a text sequence using an XML-inspired prompt structure. This formatting approach
allows for consistent parsing, enhances controllability during generation, and supports
both conditional and unconditional sampling in downstream experiments. The details of
data collection, game-wise contributions, and the exact schema used for formatting are
described in the following subsections.

\subsection{Source and Specification of Dataset}

The QUEST dataset is constructed by aggregating quest data from a curated collection
of eight popular role-playing games. These include titles from both Western and sandbox
RPG traditions, ensuring diversity in quest structure, theme, and complexity. The
primary source of this dataset is the open-access quest corpus published by V{\"a}rtinen et
al.~\cite{vartinen2022generating}, which includes quests from six mainstream RPGs. To enhance the dataset's
size and narrative richness, additional quests from \textit{Fallout}~\cite{fallout1} and \textit{Fallout 2}~\cite{fallout2} were
manually collected and structured using a semi-automated method involving an LLM
assistant.

\begin{table}[H]
  \centering
  \scriptsize
  \renewcommand{\arraystretch}{1.3}
  \begin{tabularx}{0.95\textwidth}{
    >{\raggedright\arraybackslash}p{6cm}
    >{\raggedright\arraybackslash}X
    >{\centering\arraybackslash}X
    >{\centering\arraybackslash}X
  }
    \toprule
    \textbf{Game Title} & \textbf{Source} & \textbf{Number of Quests} & \textbf{Percentage} \\
    \midrule
    Baldur's Gate~\cite{baldursgate} & V{\"a}rtinen et al.~\cite{vartinen2022generating} & 100 & 8.8183\% \\
    Baldur's Gate II: Shadows of Amn~\cite{baldursgate2shadowsofamn} & V{\"a}rtinen et al.~\cite{vartinen2022generating} & 94 & 8.2892\% \\
    The Elder Scrolls IV: Oblivion~\cite{theelderscrollsivoblivion} & V{\"a}rtinen et al.~\cite{vartinen2022generating} & 215 & 18.9594\% \\
    The Elder Scrolls V: Skyrim~\cite{theelderscrollsvskyrim} & V{\"a}rtinen et al.~\cite{vartinen2022generating} & 389 & 34.3034\% \\
    Minecraft~\cite{minecraft} & V{\"a}rtinen et al.~\cite{vartinen2022generating} & 100 & 8.8183\% \\
    Torchlight II~\cite{torchlightii} & V{\"a}rtinen et al.~\cite{vartinen2022generating} & 80 & 7.0547\% \\
    Fallout~\cite{fallout1} & Nukapedia~\cite{fallout1quests} & 58 & 5.1146\% \\
    Fallout 2~\cite{fallout2} & Nukapedia~\cite{fallout2quests} & 98 & 8.6420\% \\
    \midrule
    Total & \empty & 1134 & 99.9999\% \\
    \bottomrule
  \end{tabularx}
  \caption{Quest Distribution by Game Title}
\end{table}

Each quest entry is composed of a set of core and optional fields that reflect the semantic
elements typically found in narrative design: goals, tasks, characters, environments,
and consequences. The data schema was extended slightly from the original dataset to
accommodate additional information present in the manually added entries (e.g., character
motivations, groups, and auxiliary items). These data fields are clearly defined in
Appendix~\ref{appendix:quest-dataset}.

The inclusion of quests from \textit{Fallout}~\cite{fallout1} and \textit{Fallout 2}~\cite{fallout2} expanded the original dataset
by 156 unique entries, enhancing both the size and variety of the corpus. Among all the titles
represented, \textit{The Elder Scrolls V: Skyrim}~\cite{theelderscrollsvskyrim} contributes the
largest share of quests, accounting for approximately 34.30\% of the total, whereas \textit{Fallout}~\cite{fallout1} contributes the fewest, at 5.11\%.

To support richer expression during generation, the dataset retains several sparsely
populated or “rare” columns—fields that are not consistently filled across all game sources
but are preserved for flexibility in representing optional or game-specific quest attributes.
As a result, the compiled dataset offers a comprehensive and heterogeneous benchmark
for evaluating the narrative coherence, structural consistency, and alignment with game-world
logic in procedurally generated quests.

\begin{figure}[H]
  \begin{subfigure}[H]{0.3\textwidth}
    \centering
    \includegraphics[width=\textwidth,height=7cm]{fallout.jpg}
  \end{subfigure}
  \hfil
  \begin{subfigure}[H]{0.3\textwidth}
    \centering
    \includegraphics[width=\textwidth,height=7cm]{baldurs_gate_ii.png}
  \end{subfigure}
  \hfil
  \centering
  \begin{subfigure}[H]{0.3\textwidth}
    \centering
    \includegraphics[width=\textwidth,height=7cm]{skyrim.png}
  \end{subfigure}
  \caption{Cover arts for selected RPGs used in the QUEST Dataset~\cite{fallout1,baldursgate2shadowsofamn,theelderscrollsvskyrim}}
\end{figure}

\subsection{Tagging and Formatting}

To enable the dataset's compatibility with autoregressive LLMs, each quest entry is serialized
into a structured, tag-based format resembling XML. This formatting step ensures
consistency, simplifies prompt construction, and aligns with the model's tokenization
strategies during both training and inference.

The formatting approach employs custom tags that explicitly delineate each quest
field. Mandatory components—such as \texttt{<|begin\_objective|>}, \texttt{<|begin\_tasks|>}, and \texttt{<|begin\_quest\_giver|>}—are always included to guarantee minimum structure, while
optional fields are conditionally inserted based on content availability. This design accommodates
the dataset's heterogeneous nature and supports the generation of partial or
context-dependent quests.

\begin{lstlisting}[
  language=xml,
  caption={Snippet of the XML-like quest formatting template},
  xleftmargin=5.5cm,
]
<|begin_quest|>
<|begin_objective|>
[Objective goes here]
<|end_objective|>
<|begin_tasks|>
[Task 1]
[Task 2]
<|end_tasks|>
<|begin_quest_giver|>
[Name]: [Description]
<|end_quest_giver|>
...
<|end_quest|>
\end{lstlisting}

To support this transformation, a custom Python function (\texttt{format\_quest\_as\_xml})
was developed. It parses each quest dictionary, extracts structured elements (like \texttt{rewards},
\texttt{locations}, or \texttt{enemies}), and generates a well-formatted string using a consistent template.
Nested and list-based fields are joined with line breaks for model readability and
separation of elements. The complete tagging and formatting script is available in the
project repository\footnote{\url{https://github.com/DracoY-code/quest-generation}}.

This structured formatting approach plays a critical role in preparing the dataset for
model training. It enforces uniform tag usage, ensuring that all quest entries—regardless
of their original source—follow a standardized input format. The design also offers flexibility,
allowing non-essential fields to be selectively included or omitted without affecting
structural integrity.

Additionally, it facilitates controlled prompt construction, enabling fine-tuned exposure
of specific quest attributes during training and evaluation. By converting diverse,
game-specific data into a consistent and model-compatible format, this step serves as a
vital bridge between raw input and effective language model utilization. Further specifications
for the formatting procedure can be seen in Appendix~\ref{appendix:quest-dataset}.

\section{Model Configuration and Training}

This section describes the architecture, configuration, and training procedures used to
fine-tune language models for the task of procedural quest generation. Given the size and
complexity of modern LLMs, this phase incorporates efficient training strategies such as
quantization and Low-Rank Adaptation (LoRA)~\cite{hu2022lora} to make the process more computationally
tractable and memory efficient.

\subsection{Quantization Workflow}

To reduce memory usage and accelerate training, model weights were quantized using
the BitsAndBytes\footnote{\url{https://github.com/bitsandbytes-foundation/bitsandbytes}} library. This allowed loading larger models into memory-constrained
environments while preserving inference quality. The quantization configuration is based
on the 4-bit NormalFloat (NF4) format with double quantization and computation performed
in 16-bit precision (bfloat16 or float16).

\begin{lstlisting}[
  language=python,
  caption={Quantization configuration using BitsAndBytes with 4-bit NF4 precision},
  xleftmargin=3cm,
]
bnb_config = BitsAndBytesConfig(
    load_in_4bit=True,
    bnb_4bit_quant_type="nf4",
    bnb_4bit_use_double_quant=True,
    bnb_4bit_compute_dtype=torch.float16,
)
\end{lstlisting}

This setup significantly reduces the GPU VRAM footprint without degrading downstream
performance, enabling multi-billion parameter models to be trained on consumer-grade
hardware or limited cloud instances.

\begin{table}[H]
  \centering
  \scriptsize
  \renewcommand{\arraystretch}{1.3}
  \begin{tabularx}{0.95\textwidth}{
    >{\raggedright\arraybackslash}p{5cm}
    >{\centering\arraybackslash}X
  }
    \toprule
    \textbf{Parameter} & \textbf{Value} \\
    \midrule
    \textbf{Quantization Type} & 4-bit NormalFloat (NF4) \\
    \textbf{Compute Data Type} & \texttt{bfloat16} / \texttt{float16} \\
    \textbf{Double Quantization} & Enabled \\
    \textbf{Library Used} & \texttt{bitsandbytes} with Hugging Face Transformers \\
    \bottomrule
  \end{tabularx}
  \caption{Quantization Configuration Summary}
\end{table}

\subsection{LoRA Adapter Integration}

To further reduce the number of trainable parameters, Low-Rank Adaptation (LoRA)~\cite{hu2022lora}
modules were employed. LoRA injects trainable low-rank matrices into the attention
layers of transformer models, allowing efficient adaptation without modifying the base
model weights.

LoRA was selected for its efficiency and flexibility in fine-tuning LLMs. It enables
parameter-efficient training by updating only a small subset of parameters—typically less
than 1\% of the model's total—thereby significantly reducing both computational and
storage overhead. Its modular nature allows adapters to be added or removed without
altering the base model weights, facilitating plug-and-play experimentation across
different configurations.

\begin{lstlisting}[
  language=python,
  caption={LoRA configuration for causal language modeling using PEFT},
  xleftmargin=2.5cm,
]
lora_config = LoraConfig(
    r=4,
    lora_alpha=8,
    lora_dropout=0.0,
    target_modules=["q_proj", "v_proj", "o_proj"],
    bias="none",
    task_type="CAUSAL_LM",
)
\end{lstlisting}

Moreover, LoRA adapters support better generalization by mitigating
catastrophic forgetting, preserving the base model's pre-trained knowledge while
enabling effective learning of task-specific objectives. This configuration is implemented
using Hugging Face's peft\footnote{\url{https://github.com/huggingface/peft}} library, which provides seamless integration of LoRA within
transformer-based models.

\begin{table}[H]
  \centering
  \scriptsize
  \renewcommand{\arraystretch}{1.3}
  \begin{tabularx}{0.95\textwidth}{
    >{\raggedright\arraybackslash}p{5cm}
    >{\centering\arraybackslash}X
  }
    \toprule
    \textbf{Parameter} & \textbf{Value} \\
    \midrule
    Target Modules & Attention projection layers (\texttt{q\_proj}, \texttt{v\_proj}, \texttt{o\_proj}) \\
    Rank ($r$) & 4 \\
    Alpha ($\alpha$) & 8 \\
    Dropout & 0.0 \\
    Bias & None \\
    Library Used & \texttt{PEFT} (Parameter-Efficient Fine-Tuning) \\
    \bottomrule
  \end{tabularx}
  \caption{LoRA Adapter Integration Details}
\end{table}

\subsection{Training Pipeline Setup}

This training pipeline uses Hugging Face's \texttt{transformers}\footnote{\url{https://github.com/huggingface/transformers}}, \texttt{accelerate}\footnote{\url{https://github.com/huggingface/accelerate}}, and \texttt{datasets}\footnote{\url{https://github.com/huggingface/datasets}} frameworks.
The tokenized dataset is streamed into the model in batched form with support for gradient
accumulation, mixed precision, and distributed training.

\begin{table}[H]
  \centering
  \scriptsize
  \renewcommand{\arraystretch}{1.3}
  \begin{tabularx}{0.95\textwidth}{
    >{\raggedright\arraybackslash}p{5cm}
    >{\centering\arraybackslash}X
  }
    \toprule
    \textbf{Feature} & \textbf{Details} \\
    \midrule
    Frameworks & \texttt{transformers}, \texttt{datasets}, \texttt{accelerate}, \texttt{peft} \\
    Trainer & \texttt{transformers.Trainer} API \\
    Tokenization & Custom tokenizer with special XML-like tokens (e.g., \texttt{<|begin\_objective|>}) \\
    Loss function & Cross-entropy \\
    Evaluation & Perplexity, BLEU, METEOR, BERTScore on a held-out test split \\
    Logging & Optional integration with TensorBoard \\
    \bottomrule
  \end{tabularx}
  \caption{Training Pipeline Features}
\end{table}

This training loop supports both standalone LLM fine-tuning and adapter-based modular
fine-tuning and can be reused across multiple datasets or model architectures with
minimal reconfiguration.

\section{Evaluation Strategy}

To assess the quality of generated quests, a combination of quantitative and qualitative
evaluation methods is employed. These include standard language modeling metrics like
perplexity, reference-based n-gram overlap metrics such as BLEU and METEOR, semantic
similarity scoring through BERTScore, and a qualitative analysis of sample outputs to
assess structural and narrative coherence. Each metric offers complementary insight into
different aspects of the generation quality: fluency, relevance, semantic fidelity, and task-specific
performance.

\subsection{Perplexity}

Perplexity\footnote{\url{https://huggingface.co/spaces/evaluate-metric/perplexity}} is a widely used intrinsic metric for evaluating language models. It quantifies
how well a probability model predicts a given sequence. Formally, the perplexity of a
language model \( M \) on a sequence of tokens \( w_1, w_2, ..., w_N \) is defined as:

\begin{equation}
  \text{Perplexity}(M) = \exp\left(-\frac{1}{N} \sum_{i=1}^{N} \log P(w_i \mid w_{<i})\right)
\end{equation}

Lower perplexity indicates that the model assigns higher probability to the actual
next tokens in the sequence, suggesting better fluency and alignment with the target
distribution.

\subsection{BLEU}

The Bilingual Evaluation Understudy\footnote{\url{https://huggingface.co/spaces/evaluate-metric/bleu}} (BLEU) score is a reference-based metric originally
developed for machine translation. It measures the n-gram overlap between the generated
output and one or more reference texts. The BLEU score is calculated using precision
over n-grams and includes a brevity penalty to penalize overly short outputs. For up to
4-grams, the score is given by:

\begin{equation}
  \text{BLEU} = \text{BP} \cdot \exp\left( \sum_{n=1}^{4} w_n \log p_n \right)
\end{equation}

where \( p_n \) is the modified precision for n-grams, \( w_n \) is the weight (typically \( \frac{1}{4} \)), and BP
is the brevity penalty, defined as:

\begin{equation}
  \text{BP} = 
    \begin{cases}
    1 & \text{if } c > r \\
    \exp\left(1 - \frac{r}{c}\right) & \text{if } c \leq r
    \end{cases}
\end{equation}

where \( c \) is the length of the candidate and \( r \) is the length of the reference.

\subsection{METEOR}

METEOR\footnote{\url{https://huggingface.co/spaces/evaluate-metric/meteor}} (Metric for Evaluation of Translation with Explicit ORdering) is another
reference-based metric that improves upon BLEU by incorporating stemming, synonymy
matching, and flexible alignment. It computes a harmonic mean of unigram precision and
recall, with a higher weight on recall. It also applies a fragmentation penalty to discourage
disordered alignments. The overall score is:

\begin{equation}
  \text{METEOR} = F_{\text{mean}} \cdot (1 - Penalty)
\end{equation}

and

\begin{equation}
  F_{\text{mean}} = \frac{10 \cdot P \cdot R}{R + 9P}
\end{equation}

where \( P \) and \( R \) are unigram precision and recall, respectively.

\subsection{BERTScore}

BERTScore\footnote{\url{https://huggingface.co/spaces/evaluate-metric/bertscore}} evaluates semantic similarity by leveraging contextual embeddings from
pre-trained BERT-like models. Instead of relying on exact n-gram overlap, it compares
each token in the candidate with the most similar token in the reference using cosine
similarity. The final score is a combination of precision, recall, and F1 computed over
these similarity scores:

\begin{equation}
  \text{BERTScore}_{F1} = \frac{2 \cdot P \cdot R}{P + R}
\end{equation}

where precision \( P \) and recall \( R \) are calculated using cosine similarity between contextualized
embeddings of matched tokens.

\subsection{Qualitative Analysis}

In addition to automated evaluation, a human-based qualitative analysis is conducted on
a representative subset of the generated quests. This assessment focuses on several critical
dimensions: narrative coherence, which examines the logical flow and internal consistency
of the storyline; goal clarity, evaluating whether the objective and associated tasks are
clearly articulated; world consistency, which ensures alignment with the fictional game's
lore and setting; and language fluency, referring to grammatical correctness and natural
phrasing. These qualitative insights help contextualize the numerical results and capture
nuanced attributes of quest quality that are not fully measurable by automated metrics.

\begin{table}[H]
  \centering
  \scriptsize
  \renewcommand{\arraystretch}{1.3}
  \begin{tabularx}{0.95\textwidth}{
    >{\raggedright\arraybackslash}p{5cm}
    >{\raggedright\arraybackslash}X
  }
    \toprule
    \textbf{Criterion} & \textbf{Description} \\
    \midrule
    Narrative Coherence & Logical progression and internal consistency of the quest storyline \\
    Goal Clarity & Clear articulation of objectives and task sequences \\
    World Consistency & Adherence to the lore and world-specific constraints of the game \\
    Language Fluency & Correct grammar, phrasing, and stylistic naturalness \\
    \bottomrule
  \end{tabularx}
  \caption{Qualitative Evaluation Criteria for Generated Quests}
\end{table}

\newpage
